\subsection{\label{sec.gf.defs}Definitions}

The four examples above should have convinced you that generating functions
can be useful. Thus, it is worthwhile to put them on a rigorous footing by
first \textbf{defining} generating functions and then \textbf{justifying} the
manipulations we have been doing to them in the previous section (e.g.,
dividing them, solving quadratic equations, taking infinite sums, taking
derivatives, ...). We are next going to sketch how this can be done (see
\cite[Chapter 7 (in the 1st edition)]{Loehr-BC} and \cite[Chapter 7]{19s} for
some details).

First things first: Generating functions are not actually functions. They are
so-called \emph{formal power series} (short \emph{FPSs}). Roughly speaking, a
formal power series is a \textquotedblleft formal\textquotedblright\ infinite
sum of the form $a_{0}+a_{1}x+a_{2}x^{2}+\cdots$, where $x$ is an
\textquotedblleft indeterminate\textquotedblright\ (we shall soon see what
this all means). You cannot substitute $x=2$ into such a power series. (For
example, substituting $x=2$ into $\dfrac{1}{1-x}=1+x+x^{2}+x^{3}+\cdots$ would
lead to the absurd equality $\dfrac{1}{-1}=1+2+4+8+16+\cdots$.) The word
\textquotedblleft function\textquotedblright\ in \textquotedblleft generating
function\textquotedblright\ is somewhat of a historical artifact.

\subsubsection{\label{subsec.gf.defs.commrings}Reminder: Commutative rings}

In order to obtain a precise understanding of what FPSs are, we go back to
abstract algebra. We begin by recalling the concept of a commutative ring.
This is defined in any textbook on abstract algebra for more details, but we
recall the definition for the sake of completeness.

Informally, a \emph{commutative ring} is a set $K$ equipped with binary
operations $\oplus$, $\ominus$ and $\odot$ and elements $\mathbf{0}$ and
$\mathbf{1}$ that \textquotedblleft behave\textquotedblright\ like addition,
subtraction and multiplication (of numbers) and the numbers $0$ and $1$,
respectively. For example, they should satisfy rules like $\left(  a\oplus
b\right)  \odot c=\left(  a\odot c\right)  \oplus\left(  b\odot c\right)  $.

Formally, commutative rings are defined as follows:

\begin{definition}
\label{def.alg.commring}A \emph{commutative ring} means a set $K$ equipped
with three maps%
\begin{align*}
\oplus &  :K\times K\rightarrow K,\\
\ominus &  :K\times K\rightarrow K,\\
\odot &  :K\times K\rightarrow K
\end{align*}
and two elements $\mathbf{0}\in K$ and $\mathbf{1}\in K$ satisfying the
following axioms:

\begin{enumerate}
\item \emph{Commutativity of addition:} We have $a\oplus b=b\oplus a$ for all
$a,b\in K$.

(Here and in the following, we write the three maps $\oplus$, $\ominus$ and
$\odot$ infix -- i.e., we denote the image of a pair $\left(  a,b\right)  \in
K\times K$ under the map $\oplus$ by $a\oplus b$ rather than by $\oplus\left(
a,b\right)  $.)

\item \emph{Associativity of addition:} We have $a\oplus\left(  b\oplus
c\right)  =\left(  a\oplus b\right)  \oplus c$ for all $a,b,c\in K$.

\item \emph{Neutrality of zero:} We have $a\oplus\mathbf{0}=\mathbf{0}\oplus
a=a$ for all $a\in K$.

\item \emph{Subtraction undoes addition:} Let $a,b,c\in K$. We have $a\oplus
b=c$ if and only if $a=c\ominus b$.

\item \emph{Commutativity of multiplication:} We have $a\odot b=b\odot a$ for
all $a,b\in K$.

\item \emph{Associativity of multiplication:} We have $a\odot\left(  b\odot
c\right)  =\left(  a\odot b\right)  \odot c$ for all $a,b,c\in K$.

\item \emph{Distributivity:} We have%
\[
a\odot\left(  b\oplus c\right)  =\left(  a\odot b\right)  \oplus\left(  a\odot
c\right)  \ \ \ \ \ \ \ \ \ \ \text{and}\ \ \ \ \ \ \ \ \ \ \left(  a\oplus
b\right)  \odot c=\left(  a\odot c\right)  \oplus\left(  b\odot c\right)
\]
for all $a,b,c\in K$.

\item \emph{Neutrality of one:} We have $a\odot\mathbf{1}=\mathbf{1}\odot a=a$
for all $a\in K$.

\item \emph{Annihilation:} We have $a\odot\mathbf{0}=\mathbf{0}\odot
a=\mathbf{0}$ for all $a\in K$.
\end{enumerate}

[\textbf{Note:} Most authors do not include $\ominus$ in the definition of a
commutative ring, but instead require the existence of additive inverses for
all $a\in K$. This is equivalent to our definition, because if additive
inverses exist, then we can define $a\ominus b$ to be $a\oplus\overline{b}$
where $\overline{b}$ is the additive inverse of $b$.]

The operations $\oplus$, $\ominus$ and $\odot$ are called the \emph{addition},
the \emph{subtraction} and the \emph{multiplication} of the ring $K$. This
does not imply that they have any connection with the usual addition,
subtraction and multiplication of numbers; it merely means that they play
similar roles to the latter and behave similarly. When confusion is unlikely,
we will denote these three operations $\oplus$, $\ominus$ and $\odot$ by $+$,
$-$ and $\cdot$, respectively, and we will abbreviate $a\odot b=a\cdot b$ by
$ab$.

The elements $\mathbf{0}$ and $\mathbf{1}$ are called the \emph{zero} and the
\emph{unity} (or the \emph{one}) of the ring $K$. Again, this does not imply
that they equal the numbers $0$ and $1$, but merely that they play analogous
roles. We will simply call these elements $0$ and $1$ when confusion with the
corresponding numbers is unlikely.

We will use \emph{PEMDAS conventions} for the three operations $\oplus$,
$\ominus$ and $\odot$. These imply that the operation $\odot$ has higher
precedence than $\oplus$ and $\ominus$, while the operations $\oplus$ and
$\ominus$ are left-associative. Thus, for example, \textquotedblleft%
$ab+ac$\textquotedblright\ means $\left(  ab\right)  +\left(  ac\right)  $
(that is, $\left(  a\odot b\right)  \oplus\left(  a\odot c\right)  $).
Likewise, \textquotedblleft$a-b+c$\textquotedblright\ means $\left(
a-b\right)  +c=\left(  a\ominus b\right)  \oplus c$.
\end{definition}

Here are some examples of commutative rings:

\begin{itemize}
\item The sets $\mathbb{Z}$, $\mathbb{Q}$, $\mathbb{R}$ and $\mathbb{C}$ are
commutative rings. (Of course, the operations $\oplus$, $\ominus$ and $\odot$
of these rings are just the usual operations $+$, $-$ and $\cdot$ known from
high school.)

\item The set $\mathbb{N}$ is not a commutative ring, since it has no
subtraction. (It is, however, something called a \emph{commutative semiring}.)

\item The matrix ring $\mathbb{Q}^{m\times m}$ (this is the ring of all
$m\times m$-matrices with rational entries) is not a commutative ring for
$m>1$ (because it fails the \textquotedblleft commutativity of
multiplication\textquotedblright\ axiom). However, it satisfies all axioms
other than \textquotedblleft commutativity of multiplication\textquotedblright%
. This makes it a \emph{noncommutative ring}.

\item The set
\[
\mathbb{Z}\left[  \sqrt{5}\right]  =\left\{  a+b\sqrt{5}\ \mid\ a,b\in
\mathbb{Z}\right\}
\]
is a commutative ring with operations $+$, $-$ and $\cdot$ inherited from
$\mathbb{R}$. This is because any $a,b,c,d\in\mathbb{Z}$ satisfy%
\begin{align*}
\left(  a+b\sqrt{5}\right)  +\left(  c+d\sqrt{5}\right)   &  =\left(
a+c\right)  +\left(  b+d\right)  \sqrt{5}\in\mathbb{Z}\left[  \sqrt{5}\right]
;\\
\left(  a+b\sqrt{5}\right)  -\left(  c+d\sqrt{5}\right)   &  =\left(
a-c\right)  +\left(  b-d\right)  \sqrt{5}\in\mathbb{Z}\left[  \sqrt{5}\right]
;\\
\left(  a+b\sqrt{5}\right)  \left(  c+d\sqrt{5}\right)   &  =\left(
ac+5bd\right)  +\left(  ad+bc\right)  \sqrt{5}\in\mathbb{Z}\left[  \sqrt
{5}\right]  .
\end{align*}
This is called a \emph{subring} of $\mathbb{R}$ (i.e., a subset of
$\mathbb{R}$ that is closed under the operations $+$, $-$ and $\cdot$ and
therefore constitutes a commutative ring with these operations inherited from
$\mathbb{R}$).

\item For each $m\in\mathbb{Z}$, the set
\begin{align*}
\mathbb{Z}/m  &  =\left\{  \text{all residue classes modulo }m\right\} \\
&  =\left\{
\begin{array}
[c]{c}%
\text{equivalence classes of integers with respect to}\\
\text{the equivalence \textquotedblleft congruent modulo }%
m\text{\textquotedblright}\\
\text{(that is, \textquotedblleft differ by a multiple of }%
m\text{\textquotedblright)}%
\end{array}
\right\}
\end{align*}
is a commutative ring, with its operations defined by%
\[
\overline{a}+\overline{b}=\overline{a+b},\ \ \ \ \ \ \ \ \ \ \overline
{a}-\overline{b}=\overline{a-b},\ \ \ \ \ \ \ \ \ \ \overline{a}\cdot
\overline{b}=\overline{ab}.
\]
If $m>0$, then this ring $\mathbb{Z}/m$ is finite and has size $m$. It is also
known as $\mathbb{Z}/m\mathbb{Z}$ or $\mathbb{Z}_{m}$ (careful with the latter
notation; it can mean different things to different people). When $m$ is
prime, the ring $\mathbb{Z}/m$ is actually a finite field and is called
$\mathbb{F}_{m}$. (But, e.g., the ring $\mathbb{Z}/4$ is not a field, and not
the same as $\mathbb{F}_{4}$.)

\item In the examples we have seen so far, the elements of the commutative
ring either are numbers or (as in the case of matrices or residue classes)
consist of numbers. For a contrast, here is an example where they are sets:

For any two sets $X$ and $Y$, we define the \emph{symmetric difference}
$X\bigtriangleup Y$ of $X$ and $Y$ to be the set%
\begin{align*}
\left(  X\cup Y\right)  \setminus\left(  X\cap Y\right)   &  =\left(
X\setminus Y\right)  \cup\left(  Y\setminus X\right) \\
&  =\left\{  \text{all elements that belong to exactly one of }X\text{ and
}Y\right\}  .
\end{align*}


Fix a set $S$. Consider the power set $\mathcal{P}\left(  S\right)  $ of $S$
(that is, the set of all subsets of $S$). This power set $\mathcal{P}\left(
S\right)  $ is a commutative ring if we equip it with the operation
$\bigtriangleup$ as addition (that is, $X\oplus Y=X\bigtriangleup Y$ for any
subsets $X$ and $Y$ of $S$), with the same operation $\bigtriangleup$ as
subtraction (that is, $X\ominus Y=X\bigtriangleup Y$), and with the operation
$\cap$ as multiplication (that is, $X\odot Y=X\cap Y$), and with the elements
$\varnothing$ and $S$ as zero and unity (that is, with $\mathbf{0}%
=\varnothing$ and $\mathbf{1}=S$). Indeed, it is straightforward to see that
all the axioms in Definition \ref{def.alg.commring} hold for this ring. (For
example, distributivity holds because any three sets $A,B,C$ satisfy
$A\cap\left(  B\bigtriangleup C\right)  =\left(  A\cap B\right)
\bigtriangleup\left(  A\cap C\right)  $ and $\left(  A\bigtriangleup B\right)
\cap C=\left(  A\cap C\right)  \bigtriangleup\left(  B\cap C\right)  $.) This
is an example of a \emph{Boolean ring} (i.e., a ring in which $aa=a$ for each
element $a$ of the ring).

\item Here is another example of a semiring, which is rather useful in
combinatorics. Let $\mathbb{T}$ be the set $\mathbb{Z}\cup\left\{
-\infty\right\}  $, where $-\infty$ is just some extra symbol. Define two
operations $\oplus$ and $\odot$ on this set $\mathbb{T}$ by setting%
\[
a\oplus b=\max\left\{  a,b\right\}  \ \ \ \ \ \ \ \ \ \ \left(  \text{where
}\max\left\{  n,-\infty\right\}  :=n\text{ for each }n\in\mathbb{T}\right)
\]
and%
\[
a\odot b=a+b\ \ \ \ \ \ \ \ \ \ \left(  \text{where }n+\left(  -\infty\right)
:=\left(  -\infty\right)  +n:=-\infty\text{ for each }n\in\mathbb{T}\right)
.
\]
Then, $\mathbb{T}$ is a commutative semiring (i.e., it would satisfy
Definition \ref{def.alg.commring} if not for the lack of subtraction). It is
called \href{https://en.wikipedia.org/wiki/Tropical_semiring}{the
\emph{tropical semiring}} of $\mathbb{Z}$.
\end{itemize}

More examples can be found in algebra textbooks (or in \cite[\S 6.1]{detnotes}
or \cite[\S 5.2]{19s} or \cite[\S 2.1.2, \S 2.3.2, \S 2.3.3]{23wa}). \medskip

\textbf{Good news:} In any commutative ring $K$, the standard rules of
computation apply:

\begin{itemize}
\item You can compute finite sums (of elements of $K$) without specifying the
order of summation or the placement of parentheses. For example, for any
$a,b,c,d,e\in K$, we have%
\[
\left(  \left(  a+b\right)  +\left(  c+d\right)  \right)  +e=\left(  a+\left(
b+c\right)  \right)  +\left(  d+e\right)  ,
\]
so you can write the sum $a+b+c+d+e$ without putting parentheses around
anything. (This is called \textquotedblleft general
associativity\textquotedblright.)

Also, finite sums do not depend on the order of addends. For example, for any
$a,b,c,d,e\in K$, we have%
\[
a+b+c+d+e=d+b+a+e+c.
\]
(This is called \textquotedblleft general commutativity\textquotedblright.)

More formally: If $\left(  a_{s}\right)  _{s\in S}$ is any finite family of
elements of a commutative ring $K$ (this means that $S$ is a finite set, and
$a_{s}$ is an element of $K$ for each $s\in S$), then the finite sum%
\[
\sum\limits_{s\in S}a_{s}%
\]
is a well-defined element of $K$. Furthermore, such sums satisfy the usual
rules of sums (see \cite[\S 1.4.2]{detnotes}); for instance:

\begin{itemize}
\item If $S=X\cup Y$ and $X\cap Y=\varnothing$, then $\sum\limits_{s\in
S}a_{s}=\sum\limits_{s\in X}a_{s}+\sum\limits_{s\in Y}a_{s}$.

\item We have $\sum\limits_{s\in S}\left(  a_{s}+b_{s}\right)  =\sum
\limits_{s\in S}a_{s}+\sum\limits_{s\in S}b_{s}$.

\item If $S$ and $W$ are two finite sets, and if $f:S\rightarrow W$ is a map,
then $\sum_{s\in S}a_{s}=\sum_{w\in W}\sum_{\substack{s\in S;\\f\left(
s\right)  =w}}a_{s}$. (That is, you can subdivide a finite sum into a finite
sum of finite sums by bunching its addends arbitrarily.)
\end{itemize}

(See \cite[\S 2.14]{detnotes} for very detailed proofs\footnote{These proofs
are stated for numbers rather than elements of an arbitrary commutative ring
$K$, but the exact same reasoning works in an arbitrary ring $K$.}.
Alternatively, you can treat them as exercises on induction.)

If $S=\varnothing$, then $\sum\limits_{s\in S}a_{s}=0$ by definition. Such a
sum $\sum\limits_{s\in S}a_{s}$ with $S=\varnothing$ is called an empty sum.

\item The same holds for finite products. If $S=\varnothing$, then
$\prod\limits_{s\in S}a_{s}=1$ by definition.

\item If $a\in K$, then $-a$ denotes $0-a=\mathbf{0}-a\in K$.

\item If $n\in\mathbb{Z}$ and $a\in K$, then we can define an element $na\in
K$ by%
\[
na=%
\begin{cases}
\underbrace{a+a+\cdots+a}_{n\text{ times}}, & \text{if }n\geq0;\\
-\left(  \underbrace{a+a+\cdots+a}_{-n\text{ times}}\right)  , & \text{if
}n<0.
\end{cases}
\]
This generalizes the classical definition of multiplication (for integers) as
repeated addition.

\item If $n\in\mathbb{N}$ and $a\in K$, then we can define an element%
\[
a^{n}=\underbrace{aa\cdots a}_{n\text{ times}}\in K.
\]
In particular, $a^{0}=\underbrace{aa\cdots a}_{0\text{ times}}=1$ (since an
empty product is always $1$).

\item Standard rules hold:%
\begin{align*}
-\left(  a+b\right)   &  =\left(  -a\right)  +\left(  -b\right)
\ \ \ \ \ \ \ \ \ \ \text{for any }a,b\in K;\\
-\left(  -a\right)   &  =a\ \ \ \ \ \ \ \ \ \ \text{for any }a\in K;\\
\left(  n+m\right)  a  &  =na+ma\ \ \ \ \ \ \ \ \ \ \text{for any }a\in
K\text{ and }n,m\in\mathbb{Z};\\
\left(  nm\right)  a  &  =n\left(  ma\right)  \ \ \ \ \ \ \ \ \ \ \text{for
any }a\in K\text{ and }n,m\in\mathbb{Z};\\
a\left(  b-c\right)   &  =\left(  ab\right)  -\left(  ac\right)
\ \ \ \ \ \ \ \ \ \ \text{for any }a,b,c\in K;\\
\left(  ab\right)  ^{n}  &  =a^{n}b^{n}\ \ \ \ \ \ \ \ \ \ \text{for any
}a,b\in K\text{ and }n\in\mathbb{N};\\
a^{n+m}  &  =a^{n}a^{m}\ \ \ \ \ \ \ \ \ \ \text{for any }a\in K\text{ and
}n,m\in\mathbb{N};\\
a^{nm}  &  =\left(  a^{n}\right)  ^{m}\ \ \ \ \ \ \ \ \ \ \text{for any }a\in
K\text{ and }n,m\in\mathbb{N};\\
&  \ldots;\\
\left(  a+b\right)  ^{n}  &  =\sum_{k=0}^{n}\dbinom{n}{k}a^{k}b^{n-k}%
\ \ \ \ \ \ \ \ \ \ \text{for any }a,b\in K\text{ and }n\in\mathbb{N}.
\end{align*}
(The latter equality is known as the \emph{binomial theorem} or \emph{binomial
formula}.)
\end{itemize}

A further concept will be useful. Namely, if $K$ is a commutative ring, then
the notion of a $K$\emph{-module} is the straightforward generalization of the
concept of a $K$-vector space to cases where $K$ is not a field (but just a
commutative ring). Here is the definition of a $K$-module in detail:

\begin{definition}
\label{def.alg.module}Let $K$ be a commutative ring.

A $K$\emph{-module} means a set $M$ equipped with three maps
\begin{align*}
\oplus &  :M\times M\rightarrow M,\\
\ominus &  :M\times M\rightarrow M,\\
\rightharpoonup &  :K\times M\rightarrow M
\end{align*}
(notice that the third map has domain $K\times M$, not $M\times M$) and an
element $\overrightarrow{0}\in M$ satisfying the following axioms:

\begin{enumerate}
\item \emph{Commutativity of addition:} We have $a\oplus b=b\oplus a$ for all
$a,b\in M$.

(Here and in the following, we write the three maps $\oplus$, $\ominus$ and
$\rightharpoonup$ infix, just as for a commutative ring.)

\item \emph{Associativity of addition:} We have $a\oplus\left(  b\oplus
c\right)  =\left(  a\oplus b\right)  \oplus c$ for all $a,b,c\in M$.

\item \emph{Neutrality of zero:} We have $a\oplus\overrightarrow{0}%
=\overrightarrow{0}\oplus a=a$ for all $a\in M$.

\item \emph{Subtraction undoes addition:} Let $a,b,c\in M$. We have $a\oplus
b=c$ if and only if $a=c\ominus b$.

\item \emph{Associativity of scaling:} We have $u\rightharpoonup\left(
v\rightharpoonup a\right)  =\left(  uv\right)  \rightharpoonup a$ for all
$u,v\in K$ and $a\in M$.

\item \emph{Left distributivity:} We have $u\rightharpoonup\left(  a\oplus
b\right)  =\left(  u\rightharpoonup a\right)  \oplus\left(  u\rightharpoonup
b\right)  $ for all $u\in K$ and $a,b\in M$.

\item \emph{Right distributivity:} We have $\left(  u+v\right)
\rightharpoonup a=\left(  u\rightharpoonup a\right)  \oplus\left(
v\rightharpoonup a\right)  $ for all $u,v\in K$ and $a\in M$.

\item \emph{Neutrality of one:} We have $1\rightharpoonup a=a$ for all $a\in
M$.

\item \emph{Left annihilation:} We have $0\rightharpoonup a=\overrightarrow{0}%
$ for all $a\in M$.

\item \emph{Right annihilation:} We have $u\rightharpoonup\overrightarrow{0}%
=\overrightarrow{0}$ for all $u\in K$.
\end{enumerate}

[\textbf{Note:} Most authors do not include $\ominus$ in the definition of a
$K$-module. Instead, they may require the existence of additive inverses for
all $a\in M$. Just like for commutative rings, this is equivalent. Actually,
even the existence of additive inverses is not necessary, since the additive
inverse of an $a\in M$ can be constructed using the operation $\rightharpoonup
$ as $\left(  -1\right)  \rightharpoonup a$. Thus, you will sometimes see the
notion of a $K$-module defined entirely without any reference to subtraction
or additive inverses; it is still equivalent to ours.]

The operations $\oplus$, $\ominus$ and $\rightharpoonup$ are called the
\emph{addition}, the \emph{subtraction} and the \emph{scaling} (or the
$K$\emph{-action}) of the $K$-module $M$. When confusion is unlikely, we will
denote these three operations $\oplus$, $\ominus$ and $\rightharpoonup$ by
$+$, $-$ and $\cdot$, respectively, and we will abbreviate $a\rightharpoonup
b=a\cdot b$ by $ab$.

The element $\overrightarrow{0}$ is called the \emph{zero} (or the \emph{zero
vector}) of the $K$-module $M$. We will usually just call it $0$.

When $M$ is a $K$-module, the elements of $M$ are called \emph{vectors}, while
the elements of $K$ are called \emph{scalars}.

We will use \emph{PEMDAS conventions} for the three operations $\oplus$,
$\ominus$ and $\rightharpoonup$, with the operation $\rightharpoonup$ having
higher precedence than $\oplus$ and $\ominus$.
\end{definition}

This all having been said, we can now define formal power series.

